\documentclass[10pt,a4paper]{amsart}
\usepackage[margin=19mm]{geometry}
\usepackage{amsmath,amssymb,mathtools}
\usepackage[scr=rsfs,cal=euler]{mathalfa}
\usepackage{xfrac}
\usepackage[unicode,hidelinks,bookmarks=false]{hyperref}
\newcommand{\ie}{i.e.\ }
\newcommand{\cf}{cf.\ }
\newcommand{\resp}{respectively\ }
\newcommand{\Subsection}{\subsection}
\providecommand{\nobreakdash}{\nobreak\hbox{-}}
\setlength{\emergencystretch}{2em}
\multlinegap0pt
\usepackage{bm}
\usepackage{bbm}
\usepackage{dsfont}
\usepackage{xspace}
\usepackage{cancel}
\usepackage{mathtools}
\usepackage{amsthm}
\makeatletter
\@ifundefined{corollary}{\newtheorem{corollary}{Corollary}}{}
\@ifundefined{theorem}{\newtheorem{theorem}{Theorem}}{}
\makeatother
\title[Primes and Bivariate Polynomials without Constant Terms: A R]{Primes and Bivariate Polynomials without Constant Terms: A Recursive Algorithm}
\author{K. Lakshmanan}
\date{Source: arXiv:2405.10519v2 (CC BY 4.0)}
\begin{document}
\maketitle

\noindent\emph{Source and licence.} Primes and Bivariate Polynomials without Constant Terms: A Recursive Algorithm. Version arXiv:2405.10519v2 (CC BY 4.0), distributed under the Creative Commons Attribution 4.0 International licence, as recorded in the arXiv licence metadata for this exact version and shown on the abstract page https://arxiv.org/abs/2405.10519v2. Full rights notice: licenses/arxiv-2405.10519-CC-BY-4.0.txt.\par\smallskip

\noindent\textit{Editorial scope.} Counted unit: the Theorem whose source label is \texttt{mth} in the article (source0.tex lines 174--188, source label \texttt{mth}), delivered together with the article's own complete original proof of it (source0.tex lines 271--289, one \texttt{proof} environment of 19 lines). No mathematics is written, repaired or re-derived: statement and proof are the article's own text line for line. The proof is detached from the statement in the source: the article states the result at lines 174--188 and prints its proof at lines 271--289, and the proof environment's own header names the statement, so the association is the article's own. Required context retained uncounted: cor:m (lines 255--263). Every definition, display and label of those lines is retained unchanged. Omitted from this package and not redistributed: 1--173, 189--254, 264--270, 290--572. Nothing inside a retained excerpt is omitted or altered: every retained body is delivered exactly as the article's lines are written, so no omission receipt of any kind accompanies this package. Citations: the retained excerpts carry no citation command at all, so no bibliography entry of the article is transcribed and no reference is redistributed. Reference adaptation: none was needed and none was made; every reference inside the delivered unit resolves inside it, so no unresolved reference or citation is produced. Printed slips kept literal: no printed word, symbol, hypothesis or number of the counted statement or of its proof was altered. Layout and class: the article's own class is \texttt{article} with packages including amsmath, amssymb, amsthm, hyperref, geometry, fontenc; the wrapper is the lane's pinned \texttt{amsart} editorial wrapper at 10pt on A4, which restates the theorem-like environments the retained text needs and adds the article's own macro definitions that the retained text needs (none), with extra wrapper packages (bm, bbm, dsfont, xspace, cancel, mathtools). The delivered numbering is the unit's own. Assets: no retained span contains a figure environment, a graphics-inclusion command, a PDF-page inclusion, a table or a TikZ picture; no image file is copied into this package, the package declares no asset dependency and no image file is redistributed with it. Licence: version arXiv:2405.10519v2 of this article is distributed under the Creative Commons Attribution 4.0 International licence, as recorded in the arXiv licence metadata for this exact version and shown on the abstract page https://arxiv.org/abs/2405.10519v2. A full rights notice is delivered at licenses/arxiv-2405.10519-CC-BY-4.0.txt. Characters: every retained excerpt is pure ASCII, so the delivered engine, XeLaTeX with its default Latin Modern text font, reports no missing character.
	\begin{corollary}\label{cor:m}
		The following consequences hold:
		\begin{enumerate}
			\item If any of $f(x, y + f(x, b))$, $f(x, y)$, or $f(x, b)$ is prime, then the three values are pairwise coprime.
			\item If $f(x, y + f(a, b))$ is prime, then either $\gcd(f(x, y), f(a, b)) = 1$, or $f(x, y) = f(a, b)$, and it is prime. If $y = x$, then this further implies $\gcd(b, y) = 1$.
			\item If $f(x, b + f(x, b))$ is prime, then $f(x, b)$ must be prime.
			\item If $f(x + f(a, b), y + f(c, d))$ is prime, then either $\gcd(f(x, y), f(a, b), f(c, d)) = 1$, or all three values are equal and prime. If, in addition, $a = x$ and $b = y$, then $\gcd(x, y, c, d) = 1$.
		\end{enumerate}
	\end{corollary}

	\begin{theorem}\label{mth}
		Let $f(x, y)$ be a bivariate polynomial with non-negative coefficients and no constant term. Assume that for all $0 < x', y' \leq f(2,2)$ and for all positive integers $a, b, c, d, x, y$ such that
		\[
		x' = x + f(a, b), \quad y' = y + f(c, d),
		\]
		the following conditions hold:
		\begin{enumerate}
			\item If not all of $f(a,b)$, $f(c,d)$, and $f(x,y)$ are equal, then
			\[
			\gcd(f(a,b), f(c,d), f(x,y)) > 1.
			\]
			\item If $f(a,b) = f(c,d) = f(x,y)$, then this common value is not prime.
		\end{enumerate}
		Then $f(x, y)$ does not represent any prime.
	\end{theorem}

	\begin{proof}[Proof of Theorem~\ref{mth}]
		Assume that $f(x', y')$ is a prime value for some $x', y' > f(2, 2)$. By construction, there must exist $a, b, c, d, x, y$ such that
		\[
		x' = x + f(a, b), \quad y' = y + f(c, d).
		\]
		
		By hypothesis, for all such combinations with $0 < a, b, c, d, x, y \leq f(2, 2)$, one of the following holds:
		\begin{enumerate}
			\item The values $f(a, b), f(c, d), f(x, y)$ are not all equal, and
			\[
			\gcd(f(a, b), f(c, d), f(x, y)) > 1;
			\]
			\item The values are all equal, and that common value is not prime.
		\end{enumerate}
		
		Hence, in either case, $f(x', y')$ cannot be prime by 4 of Corollary \ref{cor:m} contradicting our assumption. Therefore, under the conditions of the theorem, the polynomial $f(x, y)$ cannot represent any prime values.
		
		The conclusion follows by noting that all combinations needed to verify this condition fall within a finite bounded range, up to $f(2, 2)$, making this check effectively computable.
	\end{proof}

\end{document}
